% Abstract
\chapter*{Abstract}
\addcontentsline{toc}{chapter}{Abstract}

This treatise presents the \textbf{Dialectical Cognition Framework (DCF)}\index{Dialectical Cognition Framework (DCF)}---a methodology for human-AI collaboration\index{Human-AI Collaboration} that emerged from extensive practical experience working with large language models. Unlike extraction-based approaches that treat AI as an answer machine, DCF positions LLMs as \emph{thinking mirrors}\index{Thinking Mirror}: collaborative partners in the architecture of thought itself.

The framework synthesizes insights from Socratic philosophy\index{Socratic Method}, cognitive science, and contemporary AI engineering to provide a structured approach for engaging with both conversational and agentic AI systems\index{Agentic AI}. Key contributions include: (1) the Thinking Mirror hypothesis for understanding human-LLM interaction; (2) a formalized five-phase Socratic prompting methodology\index{Socratic Prompting}; (3) principles for prompt chaining\index{Prompt Chaining} as cognitive architecture; (4) adaptations for the agentic era of autonomous AI systems; and (5) positioning within the contemporary AI methodology landscape including ACE-FCA, 12-Factor Agents, BMAD, and Ralph.

DCF operates at the \emph{micro level} of human-AI interaction---the cognitive strategy for how humans should think during checkpoints, reviews, and approvals. It serves as a ``cognitive operating system'' that runs on top of whatever agentic framework practitioners choose, filling a gap that existing methodologies leave implicit.

This work is intended for software engineers, technical writers, knowledge workers, and anyone seeking deeper engagement with AI beyond surface-level prompting or passive approval of autonomous agents.

\vspace{1cm}
\noindent\textbf{Keywords:} Large Language Models, Human-AI Collaboration, Socratic Method, Cognitive Science, Agentic AI, Prompt Engineering, Knowledge Architecture
